\subsection{The exact remaining tensor defect at the last degree}

\begin{proposition}[A-L16: Proved; tensor comparison and its last-degree defect]
\label{AL16}
Retain all the data and hypotheses of A-L15, and let
$N$ be any ungraded $S$-module. Put
\[
 C=M^{d_0+1}/B^{d_0+1}.
 \tag{AL.84}
\]
For $d_0\geq1$, the coefficient map $N\to M\otimes_SN$
induces
\[
 \begin{gathered}
 H^0(M\otimes_SN)=N,\qquad
 H^i(M\otimes_SN)=0\quad(1\leq i<d_0),\\
 H^{d_0}(M\otimes_SN)
       \cong\operatorname{Tor}^S_1(C,N).
 \end{gathered}
 \tag{AL.85}
\]
For $d_0=0$, the given retraction $e$ gives an actual
split exact sequence
\[
 0\longrightarrow N\longrightarrow H^0(M\otimes_SN)
 \longrightarrow\operatorname{Tor}^S_1(C,N)
 \longrightarrow0.
 \tag{AL.86}
\]
For every $d_0\geq0$ there is an explicit complex over
$\Gamma_{\mathbb Z}(T)$ satisfying the hypotheses of
A-L15 and an explicit $N$ for which this Tor group is
nonzero. Thus the last-degree term in these formulas
cannot be removed under the stated hypotheses.
\end{proposition}

\begin{proof}
The given coefficient inclusion and retraction split
$M$ as a cochain complex
\[
 M=S[0]\oplus K,
 \qquad K^0=\ker e\cong B^1,\quad K^i=M^i\ (i>0).
 \tag{AL.87}
\]
Here $d$ kills the coefficient copy of $S$, so its
restriction to $\ker e$ defines the indicated complex
without a choice of a further comparison. A-L15 proves
that $B^1,\ldots,B^{d_0+1}$ are flat. Consequently
tensoring each exact sequence (AL.81) in that range
preserves exactness. For $i<d_0$, the inclusion
$B^{i+1}\to M^{i+1}$ also remains injective after
tensoring: its quotient is $B^{i+2}$ and is flat.
These exact sequences give
\[
 H^i(K\otimes_SN)=0\quad(0\leq i<d_0),\qquad
 H^{d_0}(K\otimes_SN)
 =\ker(B^{d_0+1}\otimes_SN\to M^{d_0+1}\otimes_SN).
 \tag{AL.88}
\]
For $d_0=0$ the right-hand identity follows immediately
from $K^0=B^1$; it requires no earlier exact sequence.
In the exact sequence
\[
 0\to B^{d_0+1}\to M^{d_0+1}\to C\to0
 \tag{AL.89}
\]
the middle module is free. Its tensor--Tor exact
sequence therefore identifies the kernel in (AL.88)
with $\operatorname{Tor}^S_1(C,N)$. Applying the actual
splitting (AL.87) proves (AL.85)--(AL.86), including
the specified coefficient map.

We give the claimed example with all weights fixed.
Let $S=\Gamma_{\mathbb Z}(T)$ and take the direct sum
of $S[0]$ with the two-term cochain complex
\[
 S\langle d_0+2\rangle
      \xrightarrow{\ T\ }S\langle d_0+1\rangle,
 \quad\text{in degrees }d_0,d_0+1.
 \tag{AL.90}
\]
The displayed map sends the source generator to $T$
times the target generator and therefore preserves
weight. All free generator weights are nonnegative;
they also satisfy form degree at most weight. The
embedding $S\hookrightarrow\mathbb Q[T]$ shows that
multiplication by $T$ is injective. Hence the two-term
summand has zero cohomology in degree $d_0$, and the
coefficient copy is the whole $H^0(M)$. Projection
onto that copy gives the required retraction, also
when $d_0=0$.

Take the actual $S$-module $N=S/(T)$. It is nonzero,
since its weight-zero component is $\mathbb Z$.
After tensoring (AL.90) with $N$, its differential is
zero. The degree-$d_0$ cohomology of that summand is
therefore the nonzero module $N\langle d_0+2\rangle$.
This verifies the last assertion and the nonzero Tor
term via (AL.89). The same example after rationalization
has $S_{\mathbb Q}=\mathbb Q[T]$ and $N_{\mathbb Q}
=\mathbb Q$, with the identical zero specialized map.
The example is a cochain-module calculation; no
admissible simplicial PD structure is asserted for it.
\end{proof}
